\documentclass[10pt,a4paper]{amsart}
\usepackage[margin=19mm]{geometry}
\usepackage{amsmath,amssymb,mathtools}
\usepackage[scr=rsfs,cal=euler]{mathalfa}
\usepackage{xfrac}
\usepackage[unicode,hidelinks,bookmarks=false]{hyperref}
\newcommand{\ie}{i.e.\ }
\newcommand{\cf}{cf.\ }
\newcommand{\resp}{respectively\ }
\newcommand{\Subsection}{\subsection}
\providecommand{\nobreakdash}{\nobreak\hbox{-}}
\setlength{\emergencystretch}{2em}
\multlinegap0pt
\usepackage{bm}
\usepackage{bbm}
\usepackage{dsfont}
\usepackage{xspace}
\usepackage{cancel}
\usepackage{mathtools}
\usepackage{amsthm}
\makeatletter
\@ifundefined{theorem}{\newtheorem{theorem}{Theorem}}{}
\@ifundefined{proposition}{\newtheorem{proposition}[theorem]{Proposition}}{}
\makeatother
\DeclareMathOperator{\Ext}{Ext}
\title[Homotopy types of moment-angle complexes associated to almos]{Homotopy types of moment-angle complexes associated to almost linear resolutions}
\author{Steven Amelotte, Benjamin Briggs}
\date{Source: arXiv:2506.15457v1 (CC BY 4.0)}
\begin{document}
\maketitle

\noindent\emph{Source and licence.} Homotopy types of moment-angle complexes associated to almost linear resolutions. Version arXiv:2506.15457v1 (CC BY 4.0), distributed under the Creative Commons Attribution 4.0 International licence, as recorded in the arXiv licence metadata for this exact version and shown on the abstract page https://arxiv.org/abs/2506.15457v1. Full rights notice: licenses/arxiv-2506.15457-CC-BY-4.0.txt.\par\smallskip

\noindent\textit{Editorial scope.} Counted unit: the Proposition whose source label is \texttt{prop\_qlinstrand} in the article (source0.tex lines 905--907, source label \texttt{prop\_qlinstrand}), delivered together with the article's own complete original proof of it (source0.tex lines 909--924, one \texttt{proof} environment of 16 lines). No mathematics is written, repaired or re-derived: statement and proof are the article's own text line for line. Required context retained uncounted: none. Every definition, display and label of those lines is retained unchanged. Omitted from this package and not redistributed: 1--904, 908--908, 925--1256. Nothing inside a retained excerpt is omitted or altered: every retained body is delivered exactly as the article's lines are written, so no omission receipt of any kind accompanies this package. Citations: the retained excerpts carry no citation command at all, so no bibliography entry of the article is transcribed and no reference is redistributed. Reference adaptation: none was needed and none was made; every reference inside the delivered unit resolves inside it, so no unresolved reference or citation is produced. Printed slips kept literal: no printed word, symbol, hypothesis or number of the counted statement or of its proof was altered. Layout and class: the article's own class is \texttt{amsart} with packages including mathtools, amssymb, amsfonts, stmaryrd, amsmath, amsthm, enumerate, hyperref, color, inputenc, enumitem, xy, wrapfig, tikz; the wrapper is the lane's pinned \texttt{amsart} editorial wrapper at 10pt on A4, which restates the theorem-like environments the retained text needs and adds the article's own macro definitions that the retained text needs (\texttt{\textbackslash Ext}), with extra wrapper packages (bm, bbm, dsfont, xspace, cancel, mathtools). The delivered numbering is the unit's own. Assets: no retained span contains a figure environment, a graphics-inclusion command, a PDF-page inclusion, a table or a TikZ picture; no image file is copied into this package, the package declares no asset dependency and no image file is redistributed with it. Licence: version arXiv:2506.15457v1 of this article is distributed under the Creative Commons Attribution 4.0 International licence, as recorded in the arXiv licence metadata for this exact version and shown on the abstract page https://arxiv.org/abs/2506.15457v1. A full rights notice is delivered at licenses/arxiv-2506.15457-CC-BY-4.0.txt. Characters: every retained excerpt is pure ASCII, so the delivered engine, XeLaTeX with its default Latin Modern text font, reports no missing character.
\begin{proposition}\label{prop_qlinstrand}
    Let $S$ be a standard graded ring, and let $M$ be a finitely generated graded $S$-module. If $M$ is componentwise linear in the first $p+1$ steps, then the quasi-linear strand of $M$ contains $\Ext_S^{\leqslant p}(M,k)$.
\end{proposition}

\begin{proof}
We first assume that $M$ is generated in a single degree $d$. Let $F$ be the minimal fee resolution of $M$, so that $\Ext^i_S(M,k)={\rm Hom}_S(F_i,k)$. Assume that $F_i$ is generated in degree $d+i$ for $i< p$. We will show that this implies $\Ext_S^{i+1}(M,k)=\Ext_S^1(k,k)\cdot \Ext_S^i(M,k)$.

Since $F$ is a minimal resolution, for each minimal generator $f\in F_{i+1}$ we have $\partial(f)\neq 0$. For degree reasons, we may write $\partial(f)=\sum x_j g_j$ with $x_j\in S_1$ and $\{g_j\}$ forming a basis of $F_{i}$. Pick $j$ so that $x_j\neq 0$, and let $\theta: F_i\to k(-i-d)$ be dual to $g_j$ (so $\theta(g_j)=1$ and $\theta(g_k)=0$ when $j\neq k$). Also let $\sigma\in \Ext_S^1(k,k)\cong (S_1)^*$ be represented by a map $\sigma \colon S_1\to k$ be such that $\sigma( x_j)=1$. Calculating the Yoneda action shows that the element $\sigma\cdot\theta\in \Ext_S^{i+1}(M,k)={\rm Hom}_S(F_{i+1},k)$ satisfies  $(\sigma\cdot\theta)(f)=1$. Since $f$ was arbitrary, this shows that the elements of $\Ext_S^1(k,k)\cdot \Ext_S^i(M,k)$ do not simultaneously vanish on any minimal generator of $F_{i+1}$, and therefore they span $\Ext_S^{i+1}(M,k)={\rm Hom}_S(F_{i+1},k)$.

In case $M$ is not generated in a single degree, we proceed by induction on the number of degrees in which $M$ is generated, the base having just been established. If $d$ is the lowest degree of a minimal generator of $M$, then the exact sequence
\[
0\to M_{\langle d\rangle}\to M \to M/ M_{\langle d\rangle}\to 0
\]
Induces an exact sequence
\[
\Ext_S^i(M_{\langle d\rangle},k)\leftarrow \Ext_S^i(M,k) \leftarrow \Ext_S^i(M/ M_{\langle d\rangle},k).
\]
Note that when $i=0$ the left map is surjective. 
It follows that if, for $i<p$, the outer two terms can be generated from classes of degree zero under the action of $\Ext^1_S(k,k)$, then so can the term in them middle. Since both $M_{\langle d\rangle}$ and $M/M_{\langle d\rangle}$ are generated in fewer degrees than $M$, the induction hypothesis finishes the proof.
\end{proof}

\end{document}
